\chapter{Methodology}
\label{sec:methodology}

\section{Overview of the Concept}
\label{sec:research-design}

The methodology evaluates LLM-assisted graph recovery as a controlled
comparison between model choice and prompt strategy. Every model-strategy cell
uses the same Tesla EV feature set, dataset fingerprint, ground-truth graph,
parser, and metric suite. The independent variables are the model and the
prompting strategy; the dependent variables are graph-recovery metrics
computed from the predicted and reference adjacency matrices.

Let \(V\) be the ordered set of six telemetry features. The ground-truth graph
is denoted \(G^\ast=(V,E^\ast)\), where \(E^\ast\) contains the six directed
reference edges. For a model \(m\) and strategy \(s\), the parsed LLM output is
denoted \(\hat{G}_{m,s}=(V,\hat{E}_{m,s})\). The corresponding binary
adjacency matrices are \(A^\ast\) and \(\hat{A}_{m,s}\), where an entry of 1
at row \(i\), column \(j\) indicates a directed edge from feature \(i\) to
feature \(j\). Evaluation consists of comparing \(\hat{A}_{m,s}\) with
\(A^\ast\).

The scientific object is therefore not a free-form answer quality judgment. It
is a directed edge set. Each LLM output is reduced to relationships of the
form \texttt{source -> target}, filtered to the six allowed feature names, and
converted into a binary adjacency matrix. This makes it possible to report
precision, recall, F1 score, structural Hamming distance, false discovery
rate, Pearson correlation, Hamming distance, and Frobenius norm for every
model-strategy combination.

\summarynote{Methodology map}{The methodology begins with a normalized Tesla
EV feature view, constructs a compact statistical fingerprint, prompts local
LLMs to output directed edges, parses the edges into adjacency matrices, and
evaluates those matrices against the ground-truth DAG.}

\section{Prompt-to-Graph Pipeline}

Figure~\ref{fig:graph-recovery-pipeline} summarizes the conceptual framework.
The implementation is materialized in the notebook and result directory as
\texttt{.json}, \texttt{.npy}, \texttt{.csv}, \texttt{.png}, and
\texttt{.jsonld} artifacts.

\begin{figure}[!htbp]
  \centering
  \begingroup
  \newcommand{\pipelinebox}[1]{%
    \fcolorbox{thesisline}{thesisgray}{%
      \parbox[c][13mm][c]{0.19\textwidth}{\centering\small\bfseries #1}%
    }%
  }
  \setlength{\tabcolsep}{2.5pt}
  \begin{tabular}{@{}ccccccc@{}}
    \pipelinebox{Telemetry\\data} &
    \(\rightarrow\) &
    \pipelinebox{Prompt strategy\\layer} &
    \(\rightarrow\) &
    \pipelinebox{LLM graph\\proposal} &
    \(\rightarrow\) &
    \pipelinebox{Edge\\parser} \\
    &&&&&& \(\downarrow\) \\
    \multicolumn{3}{c}{} &&
    \pipelinebox{Graph-recovery\\metrics} &
    \(\leftarrow\) &
    \pipelinebox{Ground-truth\\comparison} \\
    &&&&&& \(\uparrow\) \\
    \multicolumn{5}{c}{} &&
    \pipelinebox{Adjacency\\matrix}
  \end{tabular}
  \endgroup
  \caption{Conceptual prompt-to-graph framework used in the thesis. Telemetry features are transformed into a prompted graph proposal, parsed into an adjacency matrix, and compared with the ground-truth DAG to compute graph-recovery metrics.}
  \label{fig:graph-recovery-pipeline}
\end{figure}

\begin{table}[!htbp]
  \centering
  \small
  \setlength{\tabcolsep}{3pt}
  \begin{tabularx}{\textwidth}{@{}L{34mm}Y Y@{}}
    \toprule
    Stage & Operation & Repository trace \\
    \midrule
    Dataset loading & Load the normalized Tesla EV CSV and keep the configured six features. & Notebook dataset configuration; local copy at \path{data/tesla.csv}. \\
    Fingerprinting & Summarize per-feature count, mean, standard deviation, minimum, median, maximum, and Spearman rank correlations. & \texttt{dataset\_fingerprint} in the supervisor notebook. \\
    Prompt construction & Combine feature descriptions, dataset fingerprint, and one prompt strategy. & \texttt{make\_base\_prompt} and \texttt{STRATEGIES}. \\
    Inference & Call local Ollama models using bounded generation settings. & \texttt{call\_ollama}; per-cell result JSON files. \\
    Parsing & Extract directed edges from the raw LLM response and normalize feature tokens. & \texttt{parse\_causal\_edges}. \\
    Matrix conversion & Convert parsed edges into a binary adjacency matrix using the fixed feature order. & \texttt{edges\_to\_adjacency}; saved \texttt{.npy} files. \\
    Evaluation & Compare predicted and ground-truth matrices using eight graph metrics. & \path{evaluation_matrix_tesla_ev.csv} and \path{evaluation_matrix_tesla_ev_long.csv}. \\
    Export & Save graph plots, edge tables, consensus edges, and JSON-LD. & \path{edges_kg.csv}, \path{edges_kg.jsonld}, and \path{consensus_tesla_ev.csv}. \\
    \bottomrule
  \end{tabularx}
  \caption{Methodological pipeline and repository traceability. Each stage is tied to an artifact or function in the supervisor-provided notebook and result directory.}
  \label{tab:graph-pipeline}
\end{table}

\section{Dataset Fingerprint}

The LLM is not prompted with all 65,261 rows. Instead, the notebook constructs
a compact textual fingerprint of the selected features. The fingerprint
contains two parts:

\begin{itemize}
  \item per-feature descriptive statistics: count, mean, standard deviation,
  minimum, median, and maximum, rounded for compactness;
  \item a Spearman rank-correlation matrix over the six selected telemetry
  features.
\end{itemize}

This design gives the model structured information about the dataset without
placing the full table into the prompt. It also keeps the input consistent
across strategies. Any observed difference between strategies is therefore
attributable to instruction structure rather than different data exposure.

The fingerprint is not a causal discovery algorithm. It summarizes marginal
distributions and rank correlations. The LLM then uses that information,
feature names, and prompt instructions to propose directed edges; direction is
evaluated only after parsing against the ground-truth DAG.

\section{Output Parsing and Matrix Construction}

The parser is intentionally tolerant of common LLM formatting variation. It
accepts ASCII arrows and related arrow glyphs, removes simple markdown or JSON
fences, attempts to read JSON objects containing edge lists, normalizes feature
tokens by case and punctuation, and handles chained expressions such as
\texttt{A -> B -> C} as adjacent directed edges. It then discards any edge
whose source or target is not one of the six allowed features.

The parser output is \(\hat{E}_{m,s}\), the predicted set of directed feature
pairs for one model-strategy cell. Duplicate edges are removed, and self-loops
are not retained. The remaining edges are converted into \(\hat{A}_{m,s}\)
using the fixed Tesla EV feature order:

\begin{center}
\small
\begin{tabular}{@{}lll@{}}
\texttt{electricenergyremaining} & \texttt{electricremaining} & \texttt{speed\_kmh} \\
\texttt{acceleration\_kmh2} & \texttt{mileage} & \texttt{distance\_m}
\end{tabular}
\end{center}

This parser-centered design matters because the evaluation is performed on
matrices, not on prose. If the model gives a relationship in an unsupported
format or uses a feature name outside the allowed set, that content cannot
contribute to a correct edge. The thesis therefore evaluates both relational
knowledge and compliance with the machine-readable output contract.

\section{Result Artifacts Used as Source of Truth}

The main methodology and results use the repository-local artifact folder:

\begin{quote}
\path{results/tesla_ev_graph_recovery}
\end{quote}

The most important files are listed in Table~\ref{tab:method-artifacts}. These
files are used instead of rerunning the experiment, and no numerical claims
are introduced beyond what they support.

\begin{table}[!htbp]
  \centering
  \small
  \setlength{\tabcolsep}{3pt}
  \begin{tabularx}{\textwidth}{@{}L{58mm}Y@{}}
    \toprule
    Artifact & Methodological role \\
    \midrule
    \path{Ground_Truth_UC_Tesla_Energy_Consumption_signed.json} & Documents the six-feature ground-truth graph and signed edge labels. \\
    \path{Ground_Truth_UC_Tesla_Energy_Consumption_adjacency_matrix.npy} & Matrix target used for metric computation. \\
    \path{tesla_ev/<model>/<strategy>.json} & Stores raw model output, parsed edges, score dictionary, latency, and \texttt{n\_runs}. \\
    \path{tesla_ev/<model>/<strategy>.npy} & Stores the predicted adjacency matrix for a model-strategy cell. \\
    \path{evaluation_matrix_tesla_ev.csv} & Wide evaluation table with one row per model-strategy cell and eight metric columns. \\
    \path{evaluation_matrix_tesla_ev_long.csv} & Long-format version with explicit dataset, model, family, and strategy metadata. \\
    \path{edges_kg.csv} and \path{edges_kg.jsonld} & Predicted edge exports for knowledge-graph style inspection. \\
    \path{consensus_tesla_ev.csv} & Edge support table aggregating votes across evaluated model-strategy cells. \\
    \bottomrule
  \end{tabularx}
  \caption{Primary artifacts used as methodological and numerical source of truth. The thesis does not report unsupported results outside these files.}
  \label{tab:method-artifacts}
\end{table}
